\documentclass[a4paper]{article}

\usepackage{anyfontsize}
\usepackage{amsmath}

\usepackage{indentfirst}
\setlength{\parindent}{0em}
\usepackage{autobreak}
\allowdisplaybreaks
\usepackage{parskip}
\usepackage{geometry}
\geometry{top=3cm,bottom=3cm,left=2cm,right=2cm}
\linespread{1.5}

\usepackage[fontset=none]{ctex}
\setCJKmainfont{Adobe Song Std}[BoldFont=Noto Sans CJK SC Medium]
\setCJKmonofont{Noto Sans Mono CJK SC}

\usepackage{newtxtext}
\usepackage{newtxmath}
\defaultfontfeatures{}
\usepackage{bm}
\renewcommand\mathbf\bm
\AtBeginDocument{
  \let\uglyepsilon\epsilon
  \renewcommand\epsilon\varepsilon
}

\begin{document}

\title{\textbf{天体光谱学作业}}
\author{1900011604\ \ 张植竣\ \ 北京大学}
\date{2022年5月1日}
\maketitle

\begin{enumerate}

    \item[1.]
    \begin{enumerate}
        \item[(a)] 漫线系：
        $4\,^2\mathrm{D}_{5/2} \rightarrow 3\,^2\mathrm{P}_{3/2}$，$4\,^2\mathrm{D}_{3/2} \rightarrow 3\,^2\mathrm{P}_{3/2}$，$4\,^2\mathrm{D}_{3/2} \rightarrow 3\,^2\mathrm{P}_{1/2}$。

        \item[(b)] 基线系：
        $5\,^2\mathrm{F}_{7/2} \rightarrow 3\,^2\mathrm{D}_{5/2}$，$5\,^2\mathrm{F}_{5/2} \rightarrow 3\,^2\mathrm{D}_{5/2}$，$5\,^2\mathrm{F}_{5/2} \rightarrow 3\,^2\mathrm{D}_{3/2}$。

        \item[(c)] $\Delta L=2$不符合选择定则$\Delta L=-1,0,1$

        \item[(d)] 主线系：
        $4\,^2\mathrm{P}_{3/2} \rightarrow 4\,^2\mathrm{S}_{1/2}$，$4\,^2\mathrm{P}_{1/2} \rightarrow 4\,^2\mathrm{S}_{1/2}$。
    \end{enumerate}

    \bigskip
    \item[2.]
    \begin{enumerate}
        \item[(a)]
        量子亏损的重要性：直观解释了穿透效应，解释了碱金属外层电子的能级偏离氢原子的情况。

        量子亏损与量子数$n$和$l$的关系：$\mu_{nl}$随$l$增大而减小，对于同一个$l$，$\mu_{nl}$与$n$之间关系不大。

        \item[(b)]
        $\mathrm{C}^{3+}$的能级表达式为：
        \[
            E_{nl} = -R_\infty \frac{\mu Z_{\mathrm{eff}}^2}{m_e(n-\mu_{nl})^2}
        \]
        其中$Z_{\mathrm{eff}}=6-3+1=4$，$\mu$为电子约化质量。

        \item[(c)]
        由$E_{nl} \propto \dfrac{1}{(n-\mu_{nl})^2}$得：
        \[
            E_{\mathrm{2s}} = 521078\,\mathrm{cm^{-1}},\quad n=2,\quad \mu_{\mathrm{2s}} = 0.163
        \]
        \[
            E_{\mathrm{3s}} = 217329\,\mathrm{cm^{-1}},\quad n=3,\quad \mu_{\mathrm{3s}} = 0.158
        \]
        则估计$\mu_{\mathrm{4s}} = 0.1575$，得到$E_{\mathrm{4s}} = 118857\,\mathrm{cm^{-1}}$
    \end{enumerate}

    \bigskip
    \item[3.]
    $^1\mathrm{D}_1$：
    $L=2,\quad J=1,\quad S=0$

    此处$L+S=2$，$|L_S|=2$，即$J$只能取2一个值，这里标记$J=1$是错误的。

    $^0\mathrm{D}_{5/2}$：
    $L=2,\quad J=\dfrac{5}{2},\quad S=-\dfrac{1}{2}$
    
    而$S$必须是大于等于0的，此处$S<0$错误。

    $^3\mathrm{P}_{3/2}$：
    $L=1,\quad J=\dfrac{3}{2},\quad S=1$

    此处$L+S=2$，$|L_S|=0$，即$J$只能取0、1、2三个值，这里标记$J=\dfrac{3}{2}$是错误的。

    \bigskip
    \item[4.]
    $\mathrm{1s^2 2p}$：有一个$l=1$，$s=\dfrac{1}{2}$的孤电子，则$L=1$，$S=\dfrac{1}{2}$，$J$可以取$\dfrac{1}{2}$或$\dfrac{3}{2}$。则可以产生的光谱项有：$\mathrm{^2P^o_{1/2}}$，$\mathrm{^2P^o_{3/2}}$

    $\mathrm{1s 2s 3s}$：有三个$l=0$，$s=\dfrac{1}{2}$的孤电子，则$L=0$，$S$可以有$\dfrac{1}{2}$、$\dfrac{1}{2}$、$\dfrac{3}{2}$三个不同的取值。则可以产生的光谱项有：$\mathrm{^2P^o}$

    $\mathrm{1s 2p 3p}$：有三个$s=\dfrac{1}{2}$的孤电子，$l$分别为：$l_1=0$，$l_2=0$，$l_3=0$，则$L$可以有$0$、$1$、$2$三个不同的取值，$S$可以有$\dfrac{1}{2}$、$\dfrac{1}{2}$、$\dfrac{3}{2}$三个不同的取值。则可以产生的光谱项有：$\mathrm{^2S}$，$\mathrm{^2S}$，$\mathrm{^4S}$，$\mathrm{^2P}$，$\mathrm{^2P}$，$\mathrm{^4P}$，$\mathrm{^2D}$，$\mathrm{^2D}$，$\mathrm{^4D}$

    跃迁$\mathrm{1s^2 2p} \rightarrow \mathrm{1s 2s 3s}$是电偶极禁戒的，因为有两个电子跃迁但是宇称发生了变化

    跃迁$\mathrm{1s^2 2p} \rightarrow \mathrm{1s 2p 3p}$是允许的

    跃迁$\mathrm{1s 2p 3p} \rightarrow \mathrm{1s 2s 3s}$是电偶极禁戒的，因为不满足Laporte规则（宇称没有改变）

    \bigskip
    \item[5.]
    只需要考虑$\mathrm{3d5p}$轨道上的孤电子，它们有$l_1=2$，$l_2=1$，$s_1=s_2=\dfrac{1}{2}$，相互组合就能产生$L=1,2,3$，$S=0,1$，则可以产生的光谱项有：$\mathrm{^1P}$，$\mathrm{^1D}$，$\mathrm{^1F}$，$\mathrm{^3P}$，$\mathrm{^3D}$，$\mathrm{^3F}$，其中能量最低的光谱项为$\mathrm{^3F}$，能级为$\mathrm{^3\mathrm{F}_2}$，$\mathrm{^3\mathrm{F}_3}$，$\mathrm{^3\mathrm{F}_4}$

    \bigskip
    \item[6.]
    $\mathrm{H}\alpha$谱线可以用Rydberg公式表示：
    \[
        \frac{1}{\lambda_{\mathrm{H}}} = R_{\mathrm{H}} \left(\frac{1}{2^2} - \frac{1}{3^2}\right)
    \]
    而$\mathrm{He}$的谱线可以表示为：
    \[
        \frac{1}{\lambda_{\mathrm{He}}} = R_{\mathrm{He}} \left(\frac{1}{n_f^2} - \frac{1}{n_i^2}\right)
    \]
    其中$R_{\mathrm{H}} = \dfrac{M_{\mathrm{H}}}{M_{\mathrm{H}} + m_e} R_\infty \approx R_\infty$，$R_{\mathrm{He}} = Z^2 \dfrac{M_{\mathrm{He}}}{M_{\mathrm{H}} + m_e} R_\infty \approx Z^2 R_\infty$，对于$\mathrm{He}$，$Z=2$，$R_{\mathrm{He}} \approx 4R_{\mathrm{H}}$

    为了使$\lambda_{\mathrm{He}} \approx \lambda_{\mathrm{H}}$，需要使：
    \[
        \frac{1}{2^2} - \frac{1}{3^2} = 4 \left(\frac{1}{4^2} - \frac{1}{6^2}\right) = 4 \left(\frac{1}{n_f^2} - \frac{1}{n_i^2}\right)
    \]
    则可以取$n_i=6$，$n_f=4$，此时波数为：
    \[
        \frac{1}{\lambda_{\mathrm{He}}} = R_{\mathrm{He}} \left(\frac{1}{4^2} - \frac{1}{6^2}\right) = \dfrac{4M_{\mathrm{He}}}{M_{\mathrm{H}} + m_e} R_\infty \left(\frac{1}{4^2} - \frac{1}{6^2}\right) = 1523920.465\,\mathrm{m^{-1}}
    \]
    或$\lambda_{\mathrm{He}}=6562.022\,\text{\AA}$，计算过程中的假设为$M_{\mathrm{H}} \gg m_e$和$M_{\mathrm{He}} \gg m_e$

\end{enumerate}


\end{document}
